% ch-a2-07.tex —— 第 2 篇（Xinwen Zhu）§4 末节 + 致谢 + 参考文献
% 底本 PDF 127–136 页 = 印刷页 114–123
\ysrc{114}
由与不取导出结构的版本类似的论证可知它是一个同构（细节见 [88]）。于是有下述交换图
\begin{center}
\begin{tikzcd}[column sep=large, row sep=large]
H_{K_0} \arrow[r, "\cong"] \arrow[d, "\text{猜想 4.11}", "\cong"']
  & C_c(G(F)\backslash G(\mathbb A)/G(\mathbb O),\Lambda) \arrow[d, "\text{猜想 4.49}", "\cong"'] \\
\End_{\Loc^{\mathrm{tame}}_0}\mathcal O_{\Loc^{\mathrm{ur}}_0} \arrow[r, "\cong"]
  & \Gamma(\Loc_{{}^cG,\PP^1},\omega_{\Loc_{{}^cG,\PP^1}}),
\end{tikzcd}
\end{center}
其中下方的同构由 (3.29) 推出。因此在这一特殊情形下，猜想 4.11 蕴含猜想 4.49。由于当 $\Lambda=\overline{\Q}_\ell$ 时猜想 4.11 成立（见注 4.12），猜想 4.49 在这一特殊情形下也成立。正如注 4.12 中已经提到的，这特别蕴含：在 $\overline{\Q}_\ell$ 上，$\Gamma(\Loc_{{}^cG,\PP^1},\omega_{\Loc_{{}^cG,\PP^1}})$ 集中在次数零（不过可以证明，层 $\omega_{\Loc_{{}^cG,\PP^1}}$ 的上同调幅度从上方无界）。

\begin{exam}{4.52}\label{thm:a2-4-52}
我们仍假设 $G$ 是约化群，但对 $v\in S$，$K_v$ 取 $G(\mathcal O_v)$ 的 Iwahori 子群。此时对 $v\in S$ 有
\[
\mathfrak A_{K_v}\cong\pi^{\mathrm{unip}}_*\mathcal O_{\Loc^{\mathrm{unip}}_{{}^cB,F_v}}\cong\pi^{\mathrm{unip}}_*\omega_{\Loc^{\mathrm{unip}}_{{}^cB,F_v}}.
\]
考虑
\[
\wt\Loc^{\mathrm{unip}}_{{}^cG,X,S}:=\Loc^{\mathrm{tame}}_{{}^cG,X,S}\times_{\prod_v\Loc^{\mathrm{tame}}_v}\prod_v\Loc^{\mathrm{unip}}_{{}^cB,F_v}.
\]
于是此情形下猜想 4.49 归结为
\[
C_c\big(\mathrm{Sht}_{\Delta(\bar\eta),K},\mathrm{Sat}(V)\big)\cong\Gamma\big(\wt\Loc^{\mathrm{unip}}_{{}^cG,X,S},(\mathcal W_FV)\otimes\omega_{\wt\Loc^{\mathrm{unip}}_{{}^cG,X,S}}\big).
\]
再次，当 $X=\PP^1$、$S=\{0,\infty\}$、$W=\mathbf 1$ 时，猜想 4.49 由猜想 4.19 推出。特别地，当 $\Lambda=\overline{\Q}_\ell$ 时它成立。细节见 [88]。
\end{exam}

为了在 Shtuka 的模空间与志村簇之间作类比，我们利用 §4.6 中猜想的范畴局部 Langlands 对应的形式体系，把上述猜想加以推广。固定一个有限的地点集 $T$。对（可能为空的）有限集 $I$，令 $\mathrm{Sht}_{(X-T)^I,T}$ 为 $X$ 上带 $I$-leg于 $X-T$、并在每个 $v\in T$ 处带额外leg的 $G$-shtuka 的模空间。我们简记 $\mathrm{Sht}_T$，而不用 $\mathrm{Sht}_{(X-T)^\emptyset,T}$。对每个 $v\in T$，选取单值化元 $\varpi_v\in\mathcal O_v$，并把 $\mathcal G_{\mathcal O_v}$ 看作 $\F_q[[\varpi_v]]$ 上的 parahoric 群概形，记作 $\mathcal G_v$。于是有局部 $\mathcal G_v$-shtuka 的模空间 (4.21)。存在自然态射
\[
\mathrm{Sht}_{(X-T)^I,T}\xrightarrow{\ \mathrm{res}\ }\prod_{v\in T}\mathrm{Sht}^{\mathrm{loc}}_v,
\]
它把 $X$ 上的整体 Shtuka 限制为在 $v\in T$ 处带leg的局部 Shtuka。一如既往，令 $\mathrm{Sht}_{\Delta(\bar\eta),T}$ 记 $\mathrm{Sht}_{(X-T)^I,T}$ 沿 $\bar\eta\to X-T\xrightarrow{\Delta}(X-T)^I$ 的基变换。

现在令 $T=S$ 为一个地点集，使得若 $v\notin S$，则 $\mathcal G(\mathcal O_v)$ 是约化群，且由 $t_v$ 决定\tn{底本此处作「is determined by $t_v$」，但 $t_v$ 在本文他处从未定义；上下文推断为某记号之误，无外部依据。}。在每个 $v\in S$ 处选取 $K_v\in\mathbf{Shv}(\mathrm{Sht}^{\mathrm{loc}}_v)$。这族层将充当在 $v\in S$ 处选定的「广义级结构」。沿 Newton 映射 $\mathrm{Nt}_v:\mathrm{Sht}^{\mathrm{loc}}_v\to\mathfrak B(G_v)$ 的固有前推，应当通过猜想 4.37 对应 $\Loc_v$ 上的一个（ind-）凝聚层 $\mathfrak A_{K_v}$。

\begin{conj}{4.53}\label{thm:a2-4-53}
对 $V\in\Rep({}^cG^I)$，有
\[
C_c\big(\mathrm{Sht}_{\Delta(\bar\eta),S},\mathrm{Sat}(V)\otimes\mathrm{res}^!(\boxtimes_{v\in S}K_v)\big)\cong\Gamma\big(\Loc_{{}^cG,F,S},(\mathcal W_FV)\otimes\mathrm{res}^!(\boxtimes_{v\in S}\mathfrak A_{K_v})\big).
\]
\end{conj}

\ysrc{115}
\begin{rem}{4.54}\label{thm:a2-4-54}
本猜想有一个更具概念性的表述：两个函子 $\prod_v\mathbf{Shv}(\mathfrak B(G_v),\Lambda)\to\IndCoh(\Loc_{{}^cG,F,S})$，一个用 Shtuka 模空间的上同调构造，另一个由猜想 4.37 得到，它们典范同构。细节见 [88]。
\end{rem}

我们在一些特殊情形下讨论这一猜想。

\begin{exam}{4.55}\label{thm:a2-4-55}
设 $K\subset\mathcal G(\mathbb O)$ 是猜想 4.49 中的级结构。假设 $S\supset S_K$。若在每个 $v\in S$ 处取 $K_v$ 为常值层沿 $[*/K_v]\to[*/\mathcal G(\mathcal O_v)]\hookrightarrow\mathrm{Sht}^{\mathrm{loc}}_v$（见 (4.24)）的前推，则猜想 4.53 还原为猜想 4.49，因为 $\mathrm{Sat}(V)\otimes\mathrm{res}^!(\boxtimes K_v)$ 正是 $\mathrm{Sat}(V)$ 沿 $\mathrm{Sht}_{\Delta(\bar\eta),K}\to\mathrm{Sht}_{\Delta(\bar\eta),S}$ 的前推，而 $\mathfrak A_{K_v}$ 应当恰是猜想 4.37 所预言的 $\mathfrak A_{K_v}$。
\end{exam}

\begin{exam}{4.56}\label{thm:a2-4-56}
保持上述情形，并特殊化到 $I=\{1\}$，于是 $V\in\Rep({}^cG)$。此外固定 $v_0\in S$。考虑下述图表
\[
\mathrm{Sht}_{X-S,S}\to\mathrm{Sht}_{X-(S-\{v_0\}),S-\{v_0\}}\hookleftarrow\mathrm{Sht}_S.
\]
对 $\mathrm{Sht}_{X-S,S}$ 上的层 $\mathrm{Sat}(V)\otimes\mathrm{res}^!(\boxtimes_{v\in S}K_v)$ 关于上述图表取邻近圈（nearby cycles），得到 $\mathrm{Sht}_S\otimes\overline{\F}_q$ 上的层 $R\Psi(\mathrm{Sat}(V)\otimes\mathrm{res}^!(\boxtimes K_v))$。已知存在 $\mathrm{Sht}^{\mathrm{loc}}_{v_0}\otimes\overline{\F}_q$ 上的层，记作 $\mathrm{Sat}(V)\star K_{v_0}$，使得
\[
R\Psi(\mathrm{Sat}(V)\otimes\mathrm{res}^!(\boxtimes K_v))\simeq\mathrm{res}^!(\boxtimes_{v\ne v_0}K_v\boxtimes(\mathrm{Sat}(V)\star K_{v_0})).
\]
此外，在猜想 4.37 之下，$\mathrm{Sat}(V)\star K_{v_0}$ 应当对应于 $(\mathcal W_{v_0}V)\otimes\mathfrak A_{K_{v_0}}$。于是猜想 4.53 预言一个典范同构
\[
\begin{aligned}
&C_c\big(\mathrm{Sht}_S\otimes\overline{\F}_q,\mathrm{res}^!(\boxtimes_{v\ne v_0}K_v\boxtimes(\mathrm{Sat}(V)\star K_{v_0}))\big)\\
&\qquad\cong\Gamma\big(\Loc_{{}^cG,F,S},\mathrm{res}^!(\boxtimes_{v\ne v_0}\mathfrak A_{K_v}\boxtimes((\mathcal W_{v_0}V)\otimes\mathfrak A_{K_{v_0}}))\big).
\end{aligned}
\]
特别地，该猜想将蕴含
\[
C_c\big(\mathrm{Sht}_{\Delta(\bar\eta),S},\mathrm{Sat}(V)\otimes\mathrm{res}^!(\boxtimes_{v\in S}K_v)\big)\cong C_c\big(\mathrm{Sht}_S\otimes\overline{\F}_q,R\Psi(\mathrm{Sat}(V)\otimes\mathrm{res}^!(\boxtimes K_v))\big).
\]
\end{exam}

\begin{exam}{4.57}\label{thm:a2-4-57}
假设 $G$ 拟分裂，并带一个固定数据（pinning）。假设 $T\subset S$ 是有限地点集，使得对 $v\in T$，$\mathcal G_v$ 是由该固定数据给出的 Iwahori 群。对每个 $v$，在 $G(\breve F_v)$ 的 Iwahori–Weyl 群 $\wt W_v$ 中选取 $w_v\in\Omega_v$（见 (4.8)），使得它们在 Kottwitz 映射下的像在 $X^{\bullet}(Z^{\Gamma_F}_{\hat G})$ 中的和为零。则族 $\{w_v\}$ 给出 $G$ 的一个内型 $G'$，以及整模型 $\mathcal G'$，使 $\mathcal G'_{\mathcal O_v}=\mathcal G_{\mathcal O_v}$ 对 $v\notin T$ 成立。我们有带leg于 $S$ 的 $\mathcal G$-Shtuka 的模空间 $\mathrm{Sht}_S$，以及带leg于 $S$ 的 $\mathcal G'$-Shtuka 的模空间 $\mathrm{Sht}'_S$。对 $v\in T$ 取 $K_v$ 为常值层沿闭嵌入
\[
L^+\mathcal G_v\cdot w_v/\mathrm{Ad}_\sigma L^+\mathcal G_v\hookrightarrow\mathrm{Sht}^{\mathrm{loc}}_v\otimes\overline{\F}_q
\]
（见 (4.25)）的前推，对 $v\in S-T$ 取 $K_v$ 为例 4.55 中与级 $\mathcal G(\mathcal O_v)$ 相伴的层。则
\[
C_c(\mathrm{Sht}_S\otimes\overline{\F}_q,\mathrm{res}^!\boxtimes K_v)=C_c(\mathrm{Sht}'_S\otimes\overline{\F}_q,\Lambda).
\]
这样我们看到，$G'$ 的自守形式空间出现在 $G$ 的 Shtuka 上同调中。可以用这一点，通过 Shtuka 模空间的上同调实现 Jacquet–Langlands 转移，从而推广 [77]。我们不在此处讨论细节，因为我们将要在志村簇的框架中表述一个猜想。
\end{exam}

\ysrc{116}
\begin{exam}{4.58}\label{thm:a2-4-58}
让我们考虑与 $G$ 相伴的 Drinfeld 模簇，它们是函数域上志村簇的类比。固定 $X$ 的一个一次地点，记作 $\infty$。为简单起见，假设 $G$ 是分裂的（带固定数据），并设 $\mathcal G$ 是 $X$ 上的群概形，满足 $\mathcal G|_{X-\{\infty\}}=G\times(X-\{\infty\})$，且 $\mathcal G_\infty$ 是 Iwahori 群概形（由固定数据决定）。

设 $V_\mu$ 是 $\hat G$ 的最高权 $\mu$ 的微表示（minuscule representation）。$V_\mu$ 的中心特征标记作 $[\mu]\in X^{\bullet}(Z_{\hat G})$。设 $w_\mu\in\Omega_\infty$（见 (4.8)）是 $G(\breve F_\infty)$ 的 Iwahori–Weyl 群中唯一使 Kottwitz 映射下的像为 $-[\mu]$ 的元素。我们在一个远离 $\infty$ 的有限集 $S_K$ 上取级结构 $K\subset\mathcal G(\mathbb O)$。于是定义与 $(G,\mu,K)$ 相伴的 Drinfeld 模簇 $\mathrm{Dr}_K(\mathcal G,\mu)$ 为：$X$ 上带leg于 $\bar\eta$、奇异性被 $V_\mu$ 界住，带leg于 $\infty$、奇异性被 $w_\mu$ 界住，并带级结构 $\Lambda$ 的 $\mathcal G$-Shtuka 的模空间。例如当 $G=\mathrm{GL}_2$ 时，$V_\mu$ 是 $\hat G=\mathrm{GL}_2$ 的对偶标准表示（此时可取 $w_\mu$ 在 $\mathrm{GL}_2(\breve F_\infty)$ 中的代表为
\[
\begin{pmatrix}0&1\\ \varpi_\infty&0\end{pmatrix},
\]
其中 $\varpi_\infty$ 是 $F_\infty$ 的单值化元），这还原为原来的 Drinfeld 模曲线。

紧支上同调 $C_c(\mathrm{Dr}_\Lambda(G,\mu),\Lambda)$\tn{底本此处下标作 $\Lambda$，而该簇在前文定义为 $\mathrm{Dr}_K(\mathcal G,\mu)$；据上下文推断为 $K$ 之误，无外部依据。} 是猜想 4.53 中所考虑上同调的特殊情形。即，令 $I=\{1\}$，$S=\{\infty\}\cup S_K$。令 $K_\infty$ 为常值层 $\Lambda$ 沿
\[
[*/\mathrm{I}_b]\cong L^+\mathcal G_\infty\cdot w_\mu/\mathrm{Ad}_\sigma L^+\mathcal G_\infty\subset\mathrm{Sht}^{\mathrm{loc}}_\infty
\]
（见 (4.25)）的前推，其余 $v\ne\infty$ 且 $v\in S$ 处的 $K_v$ 如例 4.55。则
\[
C_c\big(\mathrm{Dr}_\Lambda(G,\mu),k\big)\cong C_c\big(\mathrm{Sht}_{\Delta(\bar\eta),S},\mathrm{Sat}(V)\otimes\mathrm{res}^!(\boxtimes_{v\in S_K}K_v\boxtimes K_\infty)\big).
\]
另一方面，由猜想 4.37 我们应当有 $\mathfrak A_{K_\infty}\simeq\mathfrak A_{\mathcal G_b,\mathrm{I}_b}$。于是猜想 4.53 预言
\[
C_c\big(\mathrm{Dr}_\Lambda(G,\mu),k\big)\cong\Gamma\big(\Loc_{{}^cG,F,S},\mathcal W_{F,S}V\otimes\mathrm{res}^!(\boxtimes_{v\in S_K}\mathfrak A_{K_v}\otimes\mathfrak A_{\mathcal G_b,\mathrm{I}_b})\big).
\]
\end{exam}

\begin{exam}{4.59}\label{thm:a2-4-59}
我们也可以考虑所谓 Igusa 簇的紧支上同调。为简单起见，假设 $G$ 分裂，且 $\mathcal G=G\times X$。固定一个地点 $v_0$。令 $\mathrm{Sht}_{v_0,K}$ 为 $X$ 上带leg于 $v_0$、并在与 $v_0$ 不交的有限地点集 $S_K$ 上带 $K$-级结构的 $\mathcal G$-Shtuka 的模空间。有 $\mathrm{res}:\mathrm{Sht}_{v_0,K}\to\mathrm{Sht}^{\mathrm{loc}}_{v_0}$。设 $x$ 是 $\mathrm{Sht}^{\mathrm{loc}}_{v_0}$ 的一个 $\overline{\F}_q$-点，即在 $v_0$ 处带leg的局部 Shtuka。设 $b$ 是 $B(G_{v_0})$ 中与之相伴的元素。则自同构 $\mathrm{Aut}_x$ 是 $\overline{\F}_q$ 上的仿射群概形，且有 $[*/\mathrm{Aut}_x]\to\mathrm{Sht}^{\mathrm{loc}}_{v_0}$。$\mathrm{Sht}_{v_0,K}$ 中的中心叶（central leaf）$C_{v_0,K,x}$ 定义为纤维积
\[
C_{v_0,K,x}:=\mathrm{Sht}_{v_0,K}\times_{\mathrm{Sht}^{\mathrm{loc}}_{v_0}}[*/\mathrm{Aut}_x],
\]
而 Igusa 簇定义为纤维积
\[
\mathrm{Ig}_{v_0,K,x}:=\mathrm{Sht}_{v_0,K}\times_{\mathrm{Sht}^{\mathrm{loc}}_{v_0}}\{x\},
\]
它是 $C_{v_0,K,x}$ 上的 $\mathrm{Aut}_x$-挠子。两者的维数都是 $d=(2\rho,\nu_b)$，其中 $\nu_b$ 是 $b$ 的 Newton 点（如注 4.38）。其紧支上同调也出现在猜想 4.53 中。即，令 $I=\emptyset$，$S=\{v_0\}\cup S_K$。令 $K_{v_0}=\varinjlim_m x_{m,!}\Lambda[d]$，其中 $x_m:[*/\mathrm{Aut}_{x,m}]\to[*/\mathrm{Aut}_x]\to\mathrm{Sht}^{\mathrm{loc}}_{v_0}$，而 $\mathrm{Aut}_{x,m}\subset\mathrm{Aut}_x$ 是一族正规子群，使得 $\mathrm{Aut}_x/\mathrm{Aut}_{x,m}$（完美地）有限型。令 $K_v$（$v\in S_K$）为例 4.55 中与级 $K_v$ 相伴的层。
\end{exam}

\ysrc{117}
则
\[
C_c\big(\mathrm{Ig}_{v,K,x},\Lambda[d]\big)\cong C_c\big(\mathrm{Sht}_S,\mathrm{res}^!((\boxtimes_{v\in S_K}K_v)\otimes K_{v_0})\big).
\]
令 $\mathfrak A_{\mathcal G_b,\{1\}}$ 为注 4.38 中的 ind-凝聚层。则猜想 4.53 预言
\[
C_c\big(\mathrm{Ig}_{v,K,x},\Lambda[d]\big)\cong\Gamma\big(\Loc_{{}^cG,F,S},\mathrm{res}^!((\boxtimes_{v\in S_K}\mathfrak A_{K_v})\boxtimes\mathfrak A_{\mathcal G_b,\{1\}})\big).
\]

现在转向数域情形。事实上，通过志村簇上同调研究 Jacquet–Langlands 转移的工作 [77] 促成了这里讨论的所有猜想。因此，我们应当期待志村簇上同调也有类似的猜想性公式\footnote{探究局部对称空间的上同调是否容许类似的描述，将是很有意思的。}，尽管目前还缺少对 $\ell$ 与 $\infty$ 之上各地点处 $\mathfrak A_{K_v}$ 的描述。（特别地，$\ell$ 或 $\infty$ 处的层预期编码关于「权」的信息。）此外，数域情形下我们还没有整体 Langlands 参数的叠。因此，我们把猜想 4.49 与 4.53 的数域类比之精确表述推迟到 [26]。

这里我们表述一个猜想，它将是 [77] 主要结果之一的推广，并将蕴含在 $\{p,\infty\}$ 处不同的内型之间的 Jacquet–Langlands 对应的几何实现（工作 [77] 只给出在 $\infty$ 处不同的内型之间的 JL 转移）。设 $(G,X)$ 是志村数据。设 $V_\mu$ 记 $\hat G$ 的、与 $G$ 的志村余特征标按通常方式相伴的最高权 $\mu$ 的不可约表示。设 $p$ 是素数，$\mathcal G_p$ 是 $G_{\Q_p}$ 的 parahoric 模型。设 $K=K^pK_p$ 是级，其中 $K_p=\mathcal G_p(\Z_p)$。回忆我们（为简单起见）假设 $G$ 的中心 $Z_G$ 连通。此外作下述假设：
\begin{itemize}
\item $Z_G$ 中的极大各向异性环面在 $\R$ 上各向异性；
\item 群 $G$ 满足 Hasse 原理；
\item $h:\mathbb S\to G_\R$ 的 $G(\R)$-共轭类 $X$ 事实上是一个 $G_{\mathrm{ad}}(\R)$-共轭类。
\end{itemize}
第一条假设对于把志村簇与局部 shtuka 的模空间联系起来是本质的。后两条假设是为简化叙述而加的。若在后文中考虑志村簇的某个并，则可以去掉它们。

设 $\mathrm{Sh}_K(G,X)$ 为相应的志村簇（定义在反射域 $E$ 上），并假设它有典范的 mod $p$ 约化。令 $\mathrm{Sh}_{\mathcal G,\mu,K}$ 记 mod $p$ 纤维基变换到 $\overline{\F}_p$ 后的完美化。令 $\mathrm{Sht}^{\mathrm{loc}}_p$ 记相应的、在 $p$ 处带leg的局部 $\mathcal G_p$-shtuka 的模空间，同样基变换到 $\overline{\F}_p$。假设存在完美光滑态射
\[
\mathrm{res}:\mathrm{Sh}_{\mathcal G,\mu,K}\to\mathrm{Sht}^{\mathrm{loc}}_{p,\mu},
\]
其中 $\mathrm{Sht}^{\mathrm{loc}}_{p,\mu}\subset\mathrm{Sht}^{\mathrm{loc}}_p$ 是由那些在适当意义下奇异性被 $\mu$ 界住的局部 $\mathcal G_p$-shtuka 构成的闭子叠。我们指出，当 $(G,X)$ 是阿贝尔型时，这样的 mod $p$ 纤维 $\mathrm{Sh}_{\mathcal G,\mu,K}$ 在 [45] 中构造，而态射 $\mathrm{res}$ 在 [68] 中于一些温和限制下构造。

\ysrc{118}
现在对 $K_p\in\mathbf{Shv}(\mathrm{Sht}^{\mathrm{loc}}_{p,\mu})$，我们在 $\mathrm{Sh}_{\mathcal G,\mu,K}\otimes\overline{\F}_p$ 上得到层 $\mathrm{res}^!K_p$。如猜想 4.53 那样，可以考虑紧支上同调 $C_c(\mathrm{Sh}_{\mathcal G,\mu,K},\mathrm{res}^!K_p)$。可以记住下面两个例子。
\begin{itemize}
\item 若 $\mathrm{res}^!K_p=R\Psi$ 是移位常值层 $\Lambda[d]$ 在一般纤维 $\mathrm{Sh}_K(G,X)$ 上的邻近圈，其中 $d=\dim\mathrm{Sh}_{\mathcal G,\mu,K}$，则由 [51, 5.20]，$C_c(\mathrm{Sh}_{\mathcal G,\mu,K},\mathrm{res}^!K_p)$ 同构于 $\mathrm{Sh}_K(G,X)$ 的（移位）紧支上同调，而 $\mathfrak A_{K_p}$ 应当是 $(W_pV)\otimes\mathfrak A_{K_p}$\tn{底本此处作 $(w_pV)$，而猜想 4.60 中作 $(W_pV)$；据上下文推断此处为 $W_p$，无外部依据。}，如例 4.57。
\item 若 $\mathrm{res}^!K_p$ 是 Igusa 簇 $\mathrm{Ig}_{p,x,K}$ 上移位常值层 $\Lambda[d]$ 沿到 $\mathrm{Sh}_{\mathcal G,\mu,K}$ 的前推，其中 $d$ 是 $\mathrm{Ig}_{p,x,K}$ 的维数，则 $C_c(\mathrm{Sh}_{\mathcal G,\mu,K},\mathrm{res}^!K_p)$ 同构于 $C_c(\mathrm{Ig}_{p,x,K},\Lambda[d])$，而 $\mathfrak A_{K_p}$ 应当是 $\mathfrak A_{\mathcal G_b,\{1\}}$，如例 4.59。
\end{itemize}

现在设 $(G,X)$ 与 $(G',X')$ 是两个满足上述条件的志村数据，并对它们各自固定辅助选择。设 $p$ 是素数。假设存在内扭 $\Psi:G\to G'$（它通过 $\Psi$ 等同 $G$ 与 $G'$ 的对偶群），使得对一切 $v\ne p$ 有 $\beta_v=\beta'_v$。这特别蕴含存在良定义的同构 $\theta:G(\mathbb A^p_f)\cong G'(\mathbb A^p_f)$，在 $G(\mathbb A^p_f)$-共轭意义下唯一。固定这样的同构。设 $\mu$ 与 $\mu'$ 记相应的志村余特征标，它们给出 $\hat G$ 的不可约表示 $V_\mu$ 与 $V_{\mu'}$。

我们取与 $p$ 互素的级 $K^p\subset G(\mathbb A^p_f)$，并令 $K'^p=\theta(K^p)$。设 $K_p\subset G(\Q_p)$ 与 $K'_p\subset G'(\Q_p)$ 是 parahoric 子群。记 $H_{K^p}=H_{K'^p}$ 为相应的与 $p$ 互素的 Hecke 代数。选取 $K_p\in\mathbf{Shv}(\mathrm{Sht}^{\mathrm{loc}}_{p,\mu})$ 与 $K'_p\in\mathbf{Shv}(\mathrm{Sht}^{\mathrm{loc}}_{p,\mu'})$。猜想 4.53 提示下述猜想。

\begin{conj}{4.60}\label{thm:a2-4-60}
存在自然映射
\[
\begin{aligned}
&\Hom_{\Coh(\Loc_p)}\big((W_pV)\otimes\mathfrak A_{K_p},(W_pV')\otimes\mathfrak A_{K'_p}\big)\\
&\qquad\to\Hom_{H_{K^p}}\big(C_c(\mathrm{Sh}_{G,\mu,K},\mathrm{res}^!K_p),C_c(\mathrm{Sh}_{G',\mu',K'},\mathrm{res}^!K'_p)\big),
\end{aligned}
\]
与合成相容。当 $G=G'$、$\Psi$ 与 $\theta$ 是恒等映射、且如上 $\mathrm{res}^!K_p=\mathrm{res}^!K'_p=R\Psi$ 的特殊情形下，我们得到作用
\[
S:\End_{\Coh(\Loc_p)}\big((W_pV)\otimes\mathfrak A_{K_p}\big)\to\End_{Z^{\mathrm{tame}}_p\otimes H_{K^p}}\big(C_c(\mathrm{Sh}_K(G,X),\Lambda)\big),
\]
其中 $Z^{\mathrm{tame}}_p=H^0\Gamma(\Loc^{\mathrm{tame}}_p,\mathcal O)$ 是驯顺稳定中心 (4.13)，它应当通过映射 $Z^{\mathrm{tame}}_p\to Z(H_{K^p})$（见 (4.18)）作用在 $C_c(\mathrm{Sh}_K(G,X),\Lambda)$ 上。合成
\[
H_{K^p}\cong\End(\mathfrak A_{K_p})\to\End\big((W_pV)\otimes\mathfrak A_{K_p}\big)\xrightarrow{\ S\ }\End_{Z^{\mathrm{tame}}_p}\big(C_c(\mathrm{Sh}_K(G,X),\Lambda)\big)
\]
应当与 $H_{K^p}$ 在 $C_c(\mathrm{Sh}_K(G,X),\Lambda)$ 上的自然 Hecke 作用一致。
\end{conj}

\begin{rem}{4.61}\label{thm:a2-4-61}
工作 [77, 83] 在 $G\otimes\mathbb A_f\cong G'\otimes\mathbb A_f$ 且 $K_p$ 超特殊的情形下确认了本猜想的一个弱形式。但我们指出，即使在这一情形，猜想也更强。即，当把 $C_c(\mathrm{Sh}_K(G,X),\Lambda)$ 看作 $Z^{\mathrm{tame}}_p$-模\footnote{与 [28, 72] 中所考虑的一般局部对称空间的上同调不同，当仅把 $C_c(\mathrm{Sh}_V,k)$ 看作 $\Lambda$-模时，导出 Hecke 作用是看不见的。}时，导出 Hecke 代数 $H_{K^p}$ 作用在 $C_c(\mathrm{Sh}_K(G,X),\Lambda)$ 上。因此该猜想包含一个导出的 $S=T$ 陈述。
\end{rem}

\ysrc{119}
最后，我们简要讨论上述猜想的局部类比，即（广义）Rapoport–Zink 空间上同调的一个猜想性公式。事实上，这样的猜想性公式或多或少已包含在猜想 4.37 的等价 $\mathbb L_G$ 的猜想性质之中。

假设 $G$ 定义在局部域 $F$ 上，并设 $\mathcal G$ 是 $G$ 在 $\mathcal O$ 上的 parahoric 模型。设 $(G,b,\mu)$ 是 [62, 5.1] 意义下的局部 Shimura 数据，即 $b\in B(G)$，$\mu$ 是 $\hat G$ 的微支配权（minuscule dominant weight），且 $\kappa_G(b)=\mu|_{Z^{\Gamma_F}_{\hat G}}\in X^{\bullet}(Z^{\Gamma_F}_{\hat G})$。在此情形，Rapoport 与 Viehmann 预期存在一族刚性解析簇 $\{\mathrm{RZ}_{G,b,\mu,K}\}_\Lambda$（在 [62, §5] 中记作 $\{M^K\}$），定义在 $\breve E$ 上，以开紧子群 $K\subset\mathcal G(\mathcal O_F)$ 为指标，作为志村簇的局部类比。这里 $\breve E$ 是 $\mu$ 的反射域 $E$ 的极大非分歧扩域之完备化。对某些 $(G,b,\mu)$ 与 $K=\mathcal G(\mathcal O_F)$，$\mathrm{RZ}_{G,b,\mu,K}$ 可实现为相应 Rapoport–Zink 空间的刚性一般纤维。（这一族簇一般已在 [66, §24] 中构造。）关于这族簇的一些预期性质，我们引用 [62]；这里只提及，紧支上同调 $C_c(\mathrm{RZ}_{G,b,\mu,K}\otimes\breve E,k)$ 应当带有 $H_K\times W_E\times G_b(F)$ 的作用，且作为 $G_b(F)$-表示应属于 $\Rep_{\mathrm{f.g.}}(G_b(F),\Lambda)$。令 $\mathfrak A_{\mathcal G_b,\{1\}}$ 为注 4.38 中的 ind-凝聚层。

\begin{conj}{4.62}\label{thm:a2-4-62}
存在 $H_K\times W_E\times G_b(F)$-等变同构
\[
C_c\big(\mathrm{RZ}_{G,\mu,b,K}\otimes\breve E,\Lambda[(2\rho,\mu)]\big)\cong\Hom_{\Loc_{{}^cG,F}}\big(\mathfrak A_{\mathcal G_b,\{1\}},(\mathcal W_FV_\mu)\otimes\mathfrak A_{G,K}\big).
\]
\end{conj}
\tn{底本此处记作 $\mathrm{RZ}_{G,\mu,b,K}$，与前文的 $\mathrm{RZ}_{G,b,\mu,K}$ 次序不一致；正文照录，无外部依据。}
容易验证当 $b=1$、$\mu=0$ 时该公式成立。我们以几点注记作结。

\begin{rem}{4.63}\label{thm:a2-4-63}
(1) 首先，与整体情形类似，本猜想可看作 Kottwitz 与 Harris–Viehmann 关于 Rapoport–Zink 空间上同调的猜想（[62]）的精细化。

(2) 假设 $b$ 是基本的。可以对公式右端施用 ${}'{\mathbb D}^{Se}$，看出 $(G,\mu,b)$ 与 $(G_b,-\mu,-b)$ 的 RZ 空间的上同调在无穷远层次应当同构。这与下述事实一致：$(G,\mu,b)$ 与 $(G_b,-\mu,-b)$ 的两族簇在无穷层次同构（[62, 5.8] 与 [66, 23.3.2]）。还要注意，我们猜想 $\mathfrak A_{J_b,\{1\}}$\tn{底本此处作 $J_b$，而注 4.38 与例 4.59 之前均作 $\mathcal G_b$；据上下文推断为 $\mathcal G_b$ 之误，无外部依据。} 是连通的（ind-）凝聚层（注 4.38），而 $\mathfrak A_{G,K}$ 是通常的凝聚层（猜想 4.24(2)），所以右端只集中在非负次数。这意味着（基本的）Rapoport–Zink 空间的紧支上同调应当在中间次数以下消没，这与一般预期一致。此外，与注 4.50 类似，我们预期在 $\Loc_{{}^cG,F}$ 的孤立光滑点之上，右端只集中在次数零。

(3) 最后，把这一猜想性公式推广到非微与多leg情形（即 [66, §23] 中引入的广义 Rapoport–Zink 空间）是直接的\tn{底本此处作「is immediately.」，疑脱形容词（当作「is immediate」）；正文照译，无外部依据。}。
\end{rem}

\section*{致谢}
作者感谢 R. Bezrukavnikov、M. Emerton、T. Hemo、L. Xiao、Z. Yun 在本文写作期间的多次讨论。\ysrc{120}他感谢 M. Emerton 与 T. Feng 富有启发的讨论，这些讨论引出了猜想 4.11；并感谢 D. Ben-Zvi 关于猜想 4.19 的讨论。他感谢 P. Scholze 指出本文早期草稿中的若干不准确之处，感谢 M. Emerton 的许多宝贵意见与建议，以及 D. Hansen 的反馈。

\section*{参考文献}
\begin{refs}
\item[{[1]}] J. Ant\'onio, Moduli of $p$-adic representations of a profinite group, arXiv:1709.04275.
\item[{[2]}] D. Arinkin and D. Gaitsgory, Singular support of coherent sheaves and the geometric Langlands conjecture, Selecta Math. (N.S.) 21 (2015), no. 1, 1--199, DOI 10.1007/s00029-014-0167-5. MR3300415
\item[{[3]}] D. Arinkin, D. Gaitsgory, D. Kazhdan, S. Raskin, N. Rozenblyum, and Y. Varshavsky, The stack of local systems with restricted variation and geometric Langlands theory with nilpotent singular support, arXiv:2010.01906.
\item[{[4]}] R. Bellovin and T. Gee, $G$-valued local deformation rings and global lifts, Algebra Number Theory 13 (2019), no. 2, 333--378, DOI 10.2140/ant.2019.13.333. MR3927049
\item[{[5]}] D. Ben-Zvi, H. Chen, D. Helm, and D. Nadler, Coherent Springer theory and the categorical Deligne-Langlands correspondence, Invent. Math. 235 (2024), no. 2, 255--344, DOI 10.1007/s00222-023-01224-2. MR4689369
\item[{[6]}] R. Bezrukavnikov, On two geometric realizations of an affine Hecke algebra, Publ. Math. Inst. Hautes \'Etudes Sci. 123 (2016), 1--67, DOI 10.1007/s10240-015-0077-x. MR3502096
\item[{[7]}] G. B\"ockle, M. Harris, C. Khare, and J. A. Thorne, $\hat G$-local systems on smooth projective curves are potentially automorphic, Acta Math. 223 (2019), no. 1, 1--111, DOI 10.4310/ACTA.2019.v223.n1.a1. MR4018263
\item[{[8]}] G. B\"ockle and C. Khare, Mod $l$ representations of arithmetic fundamental groups. II. A conjecture of A. J. de Jong, Compos. Math. 142 (2006), no. 2, 271--294, DOI 10.1112/S0010437X05002022. MR2218896
\item[{[9]}] A. Borel, Automorphic $L$-functions, Automorphic forms, representations and $L$-functions (Proc. Sympos. Pure Math., Oregon State Univ., Corvallis, Ore., 1977), Proc. Sympos. Pure Math., XXXIII, Amer. Math. Soc., Providence, RI, 1979, pp. 27--61. MR546608
\item[{[10]}] A. Borel and J.-P. Serre, Sur certains sous-groupes des groupes de Lie compacts (French), Comment. Math. Helv. 27 (1953), 128--139, DOI 10.1007/BF02564557. MR54612
\item[{[11]}] A. Braverman and M. Finkelberg, A quasi-coherent description of the category $D$-mod$(Gr_{\mathrm{GL}(n)})$, Representation theory and algebraic geometry---a conference celebrating the birthdays of Sasha Beilinson and Victor Ginzburg, Trends Math., Birkh\"auser/Springer, Cham, 2022, pp. 133--149, DOI 10.1007/978-3-030-82007-7\_5. MR4486917
\item[{[12]}] K. Buzzard and T. Gee, The conjectural connections between automorphic representations and Galois representations, Proceedings of the LMS Durham Symposium, 2011. MR3444225
\item[{[13]}] W. Casselman and J. Shalika, The unramified principal series of $p$-adic groups. II. The Whittaker function, Compositio Math. 41 (1980), no. 2, 207--231. MR581582
\item[{[14]}] K. \v{C}esnavi\v{c}ius and P. Scholze, Purity for flat cohomology, Ann. of Math. (2) 199 (2024), no. 1, 51--180, DOI 10.4007/annals.2024.199.1.2. MR4681144
\item[{[15]}] T.-H. Chen and X. Zhu, Geometric Langlands in prime characteristic, Compos. Math. 153 (2017), no. 2, 395--452, DOI 10.1112/S0010437X16008113. MR3705230
\item[{[16]}] G. Chenevier, The $p$-adic analytic space of pseudocharacters of a profinite group and pseudorepresentations over arbitrary rings, Automorphic forms and Galois representations. Vol. 1, London Math. Soc. Lecture Note Ser., vol. 414, Cambridge Univ. Press, Cambridge, 2014, pp. 221--285. MR3444227
\item[{[17]}] S. Cotner, Morphisms of character varieties, Int. Math. Res. Not. IMRN 16 (2024), 11540--11548, DOI 10.1093/imrn/rnae124. MR4789091
\item[{[18]}] J.-F. Dat, Finitude pour les repr\'esentations lisses de groupes $p$-adiques (French, with English and French summaries), J. Inst. Math. Jussieu 8 (2009), no. 2, 261--333, DOI 10.1017/S1474748008000054. MR2485794
\item[{[19]}] J.-F. Dat, D. Helm, R. Kurinczuk, and G. Moss, Moduli of Langlands parameters, J. Eur. Math. Soc. (JEMS) 27 (2025), no. 5, 1827--1927, DOI 10.4171/jems/1599. MR4889237
\ysrc{121}
\item[{[20]}] J.-F. Dat, D. Helm, R. Kurinczuk, and G. Moss, Finiteness for Hecke algebras of $p$-adic groups, J. Amer. Math. Soc. 37 (2024), no. 3, 929--949, DOI 10.1090/jams/1034. MR4736530
\item[{[21]}] S. Donkin, Invariants of several matrices, Invent. Math. 110 (1992), no. 2, 389--401, DOI 10.1007/BF01231338. MR1185589
\item[{[22]}] A. J. de Jong, A conjecture on arithmetic fundamental groups, Israel J. Math. 121 (2001), 61--84, DOI 10.1007/BF02802496. MR1818381
\item[{[23]}] M. Emerton and T. Gee, Moduli stacks of \'etale $(\phi,\Gamma)$-modules and the existence of crystalline lifts, Annals of Mathematics Studies, vol. 215, Princeton University Press, Princeton, NJ, 2023, DOI 10.1515/9780691241364. MR4529886
\item[{[24]}] M. Emerton, T. Gee, and E. Hellman, An introduction to the categorical $p$-adic Langlands program, The Langlands program: II. Geometrization of the Langlands correspondence, Proc. Sympos. Pure Math., vol. 112.2, Amer. Math. Soc., Providence, RI, 2025.
\item[{[25]}] M. Emerton and D. Helm, The local Langlands correspondence for $\mathrm{GL}_n$ in families (English, with English and French summaries), Ann. Sci. \'Ec. Norm. Sup\'er. (4) 47 (2014), no. 4, 655--722, DOI 10.24033/asens.2224. MR3250061
\item[{[26]}] M. Emerton and X. Zhu, in preparation.
\item[{[27]}] L. Fargues and P. Scholze, Geometrization of the local Langlands correspondence, arXiv:2102.13459; Ast\'erisque (to appear).
\item[{[28]}] T. Feng, The spectral Hecke algebra, arXiv:1912.04413.
\item[{[29]}] V. Franjou and W. van der Kallen, Power reductivity over an arbitrary base, Doc. Math., Extra vol.: Andrei A. Suslin sixtieth birthday (2010), 171--195. MR2804253
\item[{[30]}] D. Gaitsgory, On de Jong's conjecture, Israel J. Math. 157 (2007), 155--191, DOI 10.1007/s11856-006-0006-2. MR2342444
\item[{[31]}] D. Gaitsgory, Outline of the proof of the geometric Langlands conjecture for $\mathrm{GL}_2$ (English, with English and French summaries), Ast\'erisque 370 (2015), 1--112. MR3364744
\item[{[32]}] D. Gaitsgory, From geometric to function-theoretic Langlands (or how to invent shtukas), arXiv:1606.09608.
\item[{[33]}] D. Gaitsgory, D. Kazhdan, N. Rozenblyum, and Y. Varshavsky, A toy model for the Drinfeld-Lafforgue shtuka construction, Indag. Math. (N.S.) 33 (2022), no. 1, 39--189, DOI 10.1016/j.indag.2021.11.002. MR4368480
\item[{[34]}] D. Gaitsgory and N. Rozenblyum, DG indschemes, Perspectives in representation theory, Contemp. Math., vol. 610, Amer. Math. Soc., Providence, RI, 2014, pp. 139--251, DOI 10.1090/conm/610/12080. MR3220628
\item[{[35]}] S. Galatius and A. Venkatesh, Derived Galois deformation rings, Adv. Math. 327 (2018), 470--623, DOI 10.1016/j.aim.2017.08.016. MR3762000
\item[{[36]}] A. Genestier and V. Lafforgue, Chtoucas restreints pour les groupes r\'eductifs et param\'etrisation de Langlands locale, arXiv:1709.00978.
\item[{[37]}] T. J. Haines, The stable Bernstein center and test functions for Shimura varieties, Automorphic forms and Galois representations. Vol. 2, London Math. Soc. Lecture Note Ser., vol. 415, Cambridge Univ. Press, Cambridge, 2014, pp. 118--186. MR3444233
\item[{[38]}] E. Hellmann, On the derived category of the Iwahori-Hecke algebra, Compos. Math. 159 (2023), no. 5, 1042--1110, DOI 10.1112/s0010437x23007145. MR4586564
\item[{[39]}] D. Helm, Whittaker models and the integral Bernstein center for $\mathrm{GL}_n$, Duke Math. J. 165 (2016), no. 9, 1597--1628, DOI 10.1215/00127094-3450422. MR3513570
\item[{[40]}] D. Helm, Curtis homomorphisms and the integral Bernstein center for $\mathrm{GL}_n$, Algebra Number Theory 14 (2020), no. 10, 2607--2645, DOI 10.2140/ant.2020.14.2607. MR4190413
\item[{[41]}] D. Helm and G. Moss, Converse theorems and the local Langlands correspondence in families, Invent. Math. 214 (2018), no. 2, 999--1022, DOI 10.1007/s00222-018-0816-y. MR3867634
\item[{[42]}] K. Iwasawa, On Galois groups of local fields, Trans. Amer. Math. Soc. 80 (1955), 448--469, DOI 10.2307/1992998. MR75239
\item[{[43]}] V. V. Kashin, Orbits of an adjoint and co-adjoint action of Borel subgroups of a semisimple algebraic group (Russian), Problems in group theory and homological algebra (Russian), Matematika, Yaroslav. Gos. Univ., Yaroslavl', 1990, pp. 141--158. MR1169975
\item[{[44]}] M. Kisin, Moduli of finite flat group schemes, and modularity, Ann. of Math. (2) 170 (2009), no. 3, 1085--1180, DOI 10.4007/annals.2009.170.1085. MR2600871
\item[{[45]}] M. Kisin and G. Pappas, Integral models of Shimura varieties with parahoric level structure, Publ. Math. Inst. Hautes \'Etudes Sci. 128 (2018), 121--218, DOI 10.1007/s10240-018-0100-0. MR3905466
\ysrc{122}
\item[{[46]}] R. E. Kottwitz, Isocrystals with additional structure, Compositio Math. 56 (1985), no. 2, 201--220. MR809866
\item[{[47]}] R. E. Kottwitz, Isocrystals with additional structure. II, Compositio Math. 109 (1997), no. 3, 255--339, DOI 10.1023/A:1000102604688. MR1485921
\item[{[48]}] L. Lafforgue, Chtoucas de Drinfeld et correspondance de Langlands (French, with English and French summaries), Invent. Math. 147 (2002), no. 1, 1--241, DOI 10.1007/s002220100174. MR1875184
\item[{[49]}] V. Lafforgue, Chtoucas pour les groupes r\'eductifs et param\'etrisation de Langlands globale (French), J. Amer. Math. Soc. 31 (2018), no. 3, 719--891, DOI 10.1090/jams/897. MR3787407
\item[{[50]}] V. Lafforgue and X. Zhu, D\'ecomposition au-dessus des param\`etres de Langlands elliptiques, arXiv:1811.07976.
\item[{[51]}] K.-W. Lan and B. Stroh, Nearby cycles of automorphic \'etale sheaves, Compos. Math. 154 (2018), no. 1, 80--119, DOI 10.1112/S0010437X1700745X. MR3719245
\item[{[52]}] R. P. Langlands, Representations of abelian algebraic groups, Olga Taussky-Todd: in memoriam, Pacific J. Math., Special Issue, posted on 1997, 231--250, DOI 10.2140/pjm.1997.181.231. MR1610871
\item[{[53]}] Y. F. Liu, Y. C. Tian, L. Xiao, W. Zhang, and X. W. Zhu, Deformation of rigid conjugate self-dual Galois representations, Acta Math. Sin. (Engl. Ser.) 40 (2024), no. 7, 1599--1644, DOI 10.1007/s10114-024-1409-x. MR4777059
\item[{[54]}] J. Lurie, Higher topos theory, Annals of Mathematics Studies, vol. 170, Princeton University Press, Princeton, NJ, 2009, DOI 10.1515/9781400830558. MR2522659
\item[{[55]}] J. Lurie, Higher Algebra, available at https://www.math.ias.edu/\textasciitilde lurie/.
\item[{[56]}] J. Lurie, Spectral algebraic geometry, available at https://www.math.ias.edu/\textasciitilde lurie/.
\item[{[57]}] J. Lurie, Derived algebraic geometry, available at https://www.math.ias.edu/\textasciitilde lurie/papers/DAG.pdf.
\item[{[58]}] J. Manning, Patching and multiplicity $2^k$ for Shimura curves, Algebra Number Theory 15 (2021), no. 2, 387--434, DOI 10.2140/ant.2021.15.387. MR4243652
\item[{[59]}] B. M. S. Martin, Reductive subgroups of reductive groups in nonzero characteristic, J. Algebra 262 (2003), no. 2, 265--286, DOI 10.1016/S0021-8693(03)00189-3. MR1971039
\item[{[60]}] I. Mirkovi\'c and K. Vilonen, Geometric Langlands duality and representations of algebraic groups over commutative rings, Ann. of Math. (2) 166 (2007), no. 1, 95--143, DOI 10.4007/annals.2007.166.95. MR2342692
\item[{[61]}] G. Prasad and J.-K. Yu, On finite group actions on reductive groups and buildings, Invent. Math. 147 (2002), no. 3, 545--560, DOI 10.1007/s002220100182. MR1893005
\item[{[62]}] M. Rapoport and E. Viehmann, Towards a theory of local Shimura varieties, M\"unster J. Math. 7 (2014), no. 1, 273--326. MR3271247
\item[{[63]}] F. Rodier, Mod\`ele de Whittaker et caract\`eres de repr\'esentations, in Non-commutative harmonic analysis, Lecture Notes in Math. 466 (1975), 151--171.
\item[{[64]}] P. Schneider, Smooth representations and Hecke modules in characteristic $p$, Pacific J. Math. 279 (2015), no. 1-2, 447--464, DOI 10.2140/pjm.2015.279.447. MR3437786
\item[{[65]}] P. Scholze, The local Langlands correspondence for $\mathrm{GL}_n$ over $p$-adic fields, Invent. Math. 192 (2013), no. 3, 663--715, DOI 10.1007/s00222-012-0420-5. MR3049932
\item[{[66]}] P. Scholze and J. Weinstein, Moduli of $p$-divisible groups, Camb. J. Math. 1 (2013), no. 2, 145--237, DOI 10.4310/CJM.2013.v1.n2.a1. MR3272049
\item[{[67]}] P. Scholze and J. Weinstein, Berkeley lectures on $p$-adic geometry, Annals of Mathematics Studies, vol. 207, Princeton University Press, Princeton, NJ, 2020. MR4446467
\item[{[68]}] X. Shen, C.-F. Yu, and C. Zhang, EKOR strata for Shimura varieties with parahoric level structure, Duke Math. J. 170 (2021), no. 14, 3111--3236, DOI 10.1215/00127094-2021-0047. MR4319228
\item[{[69]}] B. To\"en, Derived algebraic geometry, EMS Surv. Math. Sci. 1 (2014), no. 2, 153--240, DOI 10.4171/EMSS/4. MR3285853
\item[{[70]}] W. van der Kallen, Good Grosshans filtration in a family (English, with English and French summaries), Autour des sch\'emas en groupes. Vol. III, Panor. Synth\`eses, vol. 47, Soc. Math. France, Paris, 2015, pp. 111--129. MR3525842
\item[{[71]}] W. van der Kallen, Reductivity properties over an affine base, Indag. Math. (N.S.) 32 (2021), no. 5, 961--967, DOI 10.1016/j.indag.2020.09.009. MR4310008
\item[{[72]}] A. Venkatesh, Derived Hecke algebra and cohomology of arithmetic groups, Forum Math. Pi 7 (2019), e7, 119, DOI 10.1017/fmp.2019.6. MR4061961
\ysrc{123}
\item[{[73]}] M.-F. Vign\'eras, The pro-$p$-Iwahori Hecke algebra of a reductive $p$-adic group III (spherical Hecke algebras and supersingular modules), J. Inst. Math. Jussieu 16 (2017), no. 3, 571--608, DOI 10.1017/S1474748015000146. MR3646282
\item[{[74]}] E. B. Vinberg, On invariants of a set of matrices, J. Lie Theory 6 (1996), no. 2, 249--269. MR1424635
\item[{[75]}] C. Wang-Erickson, Algebraic families of Galois representations and potentially semi-stable pseudodeformation rings, Math. Ann. 371 (2018), no. 3-4, 1615--1681, DOI 10.1007/s00208-017-1557-8. MR3831282
\item[{[76]}] M. Weidner, Pseudocharacters of homomorphisms into classical groups, Transform. Groups 25 (2020), no. 4, 1345--1370, DOI 10.1007/s00031-020-09603-2. MR4166690
\item[{[77]}] L. Xiao and X. Zhu, Cycles on Shimura varieties via geometric Satake, arXiv:1707.05700, 2017.
\item[{[78]}] L. Xiao and X. Zhu, On vector-valued twisted conjugation invariant functions on a group, Representations of reductive groups, Proc. Sympos. Pure Math., vol. 101, Amer. Math. Soc., Providence, RI, 2019, pp. 361--425, DOI 10.1090/pspum/101/14. With an appendix by Stephen Donkin. MR3930024
\item[{[79]}] C. Xue, Finiteness of cohomology groups of stacks of shtukas as modules over Hecke algebras, and applications (English, with English and French summaries), \'Epijournal G\'eom. Alg\'ebrique 4 (2020), Art. 6, 42, DOI 10.46298/epiga.2020.volume4.5550. MR4113657
\item[{[80]}] C. Xue, Cohomology with integral coefficients of stacks of shtukas, arXiv:2001.05805v1.
\item[{[81]}] C. Xue, Smoothness of cohomology sheaves of stacks of shtukas, arXiv:2012.12833.
\item[{[82]}] J. Yu, The integral geometric Satake equivalence in mixed characteristic, Represent. Theory 26 (2022), 874--905, DOI 10.1090/ert/610. MR4470196
\item[{[83]}] J. Yu, A geometric Jacquet-Langlands transfer for automorphic forms of higher weights, Trans. Amer. Math. Soc. 375 (2022), no. 10, 6843--6873, DOI 10.1090/tran/8611. MR4491415
\item[{[84]}] X. Zhu, The geometric Satake correspondence for ramified groups (English, with English and French summaries), Ann. Sci. \'Ec. Norm. Sup\'er. (4) 48 (2015), no. 2, 409--451, DOI 10.24033/asens.2248. MR3346175
\item[{[85]}] X. Zhu, Affine Grassmannians and the geometric Satake in mixed characteristic, Ann. of Math. (2) 185 (2017), no. 2, 403--492, DOI 10.4007/annals.2017.185.2.2. MR3612002
\item[{[86]}] X. Zhu, Geometric Satake, categorical traces, and arithmetic of Shimura varieties, Current developments in mathematics 2016, Int. Press, Somerville, MA, 2018, pp. 145--206. MR3837875
\item[{[87]}] X. Zhu, A note on integral Satake isomorphisms, Arithmetic geometry, Tata Inst. Fundam. Res. Stud. Math., vol. 41, Tata Inst. Fund. Res., Mumbai, 2024, pp. 469--489. MR4812712
\item[{[88]}] X. Zhu, Tame categorical local Langlands correspondence, arXiv:2504.07482, 2025.
\end{refs}

\medskip
\noindent Department of Mathematics, Stanford University, Stanford, CA 94305

\noindent Email address: \texttt{zhuxw@stanford.edu}
